%% --------------------------------------------------------
\section{Співвідношення невизначеностей}\label{sec:uncert_principle}
%% --------------------------------------------------------

Міркування цього розділу застосовні до будь-яких процесів, що можна аналізувати за допомогою перетворення Фур’є, зокрема, до хвильових пакетів, які є
суперпозицією плоских хвиль з різними частотами.

Нехай маємо деякий процес, що можна описати за допомогою функції (взагалі кажучи, комплексної) від дійсної змінної \( t \), причому для цієї функції існує перетворення Фур’є:

\begin{equation*}
\phi(t) = \frac{1}{\sqrt{2\pi}} \int d\omega \cdot e^{-i\omega t} \tilde\phi(\omega)
\end{equation*}

Вважатимемо, що \( |\phi(t)| \) та \( |\tilde\phi(\omega)| \) досить швидко спадають, якщо відповідно \( |t| \to \infty \) та \( |\omega| \to \infty \). Також вважаємо, що середнє за ансамблем\footnote{Як і раніше, середні за ансамблем позначаємо $\langle\ldots\rangle$ (без індексу), на відміну, наприклад, від $\langle\ldots\rangle_{\phi}$. } від зазначеного поля дорівнює нулю $\langle \phi(t) \rangle = 0 $,  $\langle \tilde\phi(t) \rangle = 0 $.

Введемо величини, що характеризують тривалість процесу, що відповідає функції \( \phi(t) \), та ширину його спектру відповідно до \( \tilde\phi(\omega) \).
Середні значення фізичних величин, що відповідають для окремої реалізації обчислюватимемо з вагою \( |\phi(t)|^2 \) в $t$-просторі, а в \( \omega \)-просторі --- з вагою \( |\tilde\phi(\omega)|^2 \), причому, не
зменшуючи загальності, приймемо умову нормування:

\begin{equation*}
\int dt \cdot |\phi(t)|^2 = 1
\end{equation*}

Звідси, за рівністю Парсеваля~\eqref{eq:Parseval}, в \( \omega \)-просторі також:

\begin{equation*}
\int d\omega \cdot |\tilde\phi(\omega)|^2 = 1
\end{equation*}

Відповідно, введемо (за означенням) середнє значення деякої величини \( F(t) \) у \( t \)-просторі для реалізації $\phi$, яке позначимо:

\begin{equation*}
\langle F \rangle_{\phi} = \int dt \cdot F(t) |\phi(t)|^2,
\end{equation*}
а для \( \tilde{F}(\omega) \) в \( \omega \)-просторі:

\begin{equation*}
\langle \tilde{F} \rangle_{\phi} = \int d\omega \cdot \tilde{F}(\omega) |\phi(\omega)|^2.
\end{equation*}

Очевидно, усі ці співвідношення зберігаються при усередненні за ансамблем\footnote{Далі вважаємо, що для випадкового процесу виконано умови, що дозволяють міняти порядок операцій усереднення та інтегрування.}:
\begin{equation*}
    \langle F \rangle = \int dt \cdot F(t) \left\langle|\phi(t)|^2\right\rangle,\ \langle \tilde{F} \rangle = \int d\omega \cdot \tilde{F}(\omega) \left\langle |\phi(\omega)|^2\right\rangle.
\end{equation*}


Виберемо відлік часу таким чином, щоб:

\begin{equation*}
\langle t \rangle = \int dt \cdot t \left\langle|\phi(t)|^2\right\rangle = 0
\end{equation*}

Тоді тривалість процесу \( \phi(t) \) можна описати середньоквадратичним значенням:

\begin{equation}
\langle \Delta t^2 \rangle = \int dt \cdot t^2 \langle |\phi(t)|^2 \rangle
\label{eq:time_duration}
\end{equation}

Відповідно, в \( \omega \)-просторі середнє значення частоти є:

\begin{equation*}
\omega_a = \langle \omega \rangle = \int d\omega \cdot \omega \langle |\tilde\phi(\omega)|^2 \rangle,
\end{equation*}
а ширину спектру частот можна описати величиною:

\begin{equation}
\langle \Delta \omega^2 \rangle = \langle (\omega - \omega_a)^2 \rangle = \int d\omega \cdot (\omega - \omega_a)^2 \langle |\tilde\phi(\omega)|^2 \rangle
\label{eq:frequency_width}
\end{equation}

Покажемо, що існує спільне обмеження на величини $\langle \Delta \omega^2 \rangle$ та $\langle \Delta t^2 \rangle$. Для цього перепишемо~\eqref{eq:time_duration} через перетворення Фур’є.

Маємо:

\begin{equation*}
t \phi(t) = \frac{1}{\sqrt{2\pi}} \int d\omega \cdot i \frac{d e^{-i\omega t}}{d\omega} \tilde\phi(\omega) = \frac{1}{\sqrt{2\pi}} \int d\omega \cdot
e^{-i\omega t} \left[ -i \frac{d\tilde\phi}{d\omega} \right],
\end{equation*}
тобто помноження на \( t \) індукує диференціювання у просторі частот. Тоді формулу~\eqref{eq:time_duration} можна переписати так:

\begin{equation}
\langle \Delta t^2 \rangle =
%
\left\langle \int dt \cdot t |\phi(t)|^2 \right\rangle =
%
\left\langle
\int d\omega
\left|
\frac{d\tilde\phi}{d\omega}
\right|^2
\right\rangle
=
%
 \int d\omega
\left\langle
\left| \frac{d\tilde{\phi}}{d\omega} \right|^2
\right\rangle
%
\label{eq:time_duration_fourier}
\end{equation}
%
де застосовано рівність Парсеваля~\eqref{eq:Parseval}.

Для будь-якого дійсного \( x \):

\begin{equation*}
0 \leqslant \int d\omega \left| x \frac{d\tilde\phi}{d\omega} + (\omega - \omega_a) \tilde\phi \right|^2
\equiv
\int d\omega \left( x \frac{d\tilde\phi}{d\omega} + (\omega -
\omega_a) \tilde\phi \right) \left( x \frac{d\tilde\phi^{*}}{d\omega} + (\omega - \omega_a) \tilde\phi^{*} \right).
\end{equation*}

Звідси, після застосування усереднення за ансамблем, за допомогою~\eqref{eq:frequency_width} та~\eqref{eq:time_duration_fourier}, маємо:

\begin{equation*}
0 \leqslant  x^2 \langle \Delta t^2 \rangle
+
\langle \Delta \omega^2 \rangle
+
\left\langle x \int d\omega (\omega - \omega_a) \left[ \frac{d\tilde\phi}{d\omega}
\tilde\phi^{*} + \frac{d\tilde\phi^{*}}{d\omega} \tilde\phi \right] \right\rangle.
\end{equation*}

Вираз у квадратних дужках в останньому доданку є повною похідною, тому інтеграл у правій частині нерівності перетворимо так:

\begin{equation*}
\int d\omega (\omega - \omega_a) \frac{d}{d\omega} (\tilde\phi \tilde\phi^{*}) = -\int d\omega |\tilde\phi|^2 \frac{d}{d\omega} (\omega - \omega_a) = - \int d\omega \left| \tilde\phi\right|^2 =  -1.
\end{equation*}

Таким чином:

\begin{equation*}
x^2 \langle \Delta t^2 \rangle + \langle \Delta \omega^2 \rangle - x \geqslant0.
\end{equation*}

Це співвідношення має виконуватися для будь-яких \( x \), тому дискримінант квадратного тричлена має бути від’ємний. Звідси маємо нерівність:

\begin{equation}
\langle \Delta t^2 \rangle \langle \Delta \omega^2 \rangle \geqslant\frac{1}{4},
\label{eq:uncertainty_relation}
\end{equation}
яку називають \emph{співвідношенням невизначеностей}\index{Співвідношення невизначеностей}.

Зауважимо, що в квантовій механіці це дає зв’язок між тривалістю процесу та невизначеністю енергії. Зрозуміло, що усі міркування і співвідношення~\eqref{eq:uncertainty_relation}  можна застосовувати і до невипадкового процесу, що представлений однією реалізацією.

З~\eqref{eq:uncertainty_relation} випливає, що сигнал, що близький до монохроматичного, тобто такий, що має малий розкид частот, має бути досить
тривалим. Навпаки, короткочасний імпульсний сигнал має широкий спектр частот. Нерівність, аналогічну~\eqref{eq:uncertainty_relation}, можна отримати для
довжини хвильового пакету та ширини інтервалу хвильових векторів у певному напрямі. Квантовий відповідник ~\eqref{eq:uncertainty_relation} --- співвідношення невизначеностей для
імпульсу та координати.